% Einheit 4 von 5 – Zylinder (bank/koerper/e4.jsonl, zusammenbau v0.1)
%% TODO Zweigzeile Teil 1: was hier gelernt wird (Satz in Schülersprache aus dem Einheitstitel) – Einheit 4
\einheitenkopf[e4]{Einheit 4 von 5 · Zylinder}
\zweigzeile{ab Kl. 8 (am Gymnasium ab Kl. 7) · P10}
\setcounter{aufgabe}{21}

%% TODO Titel: Ich-kann-Satz fehlt in der Bank (Platzhalter: Kettenname) – Nr. 22
%% TODO Anweisung: ein Satz über der ersten Teilaufgabe fehlt in der Bank – Nr. 22
\begin{aufgabe}{Welche Einheit?}
\begin{teile}
\teil Gefragt ist, wie viel Blech man für den Mantel einer Dose braucht. In welcher Einheit steht das Ergebnis? \\ \kreuz{cm} \\ \kreuz{cm²} \\ \kreuz{cm³}
\end{teile}
\end{aufgabe}

%% TODO Titel: Ich-kann-Satz fehlt in der Bank (Platzhalter: Kettenname) – Nr. 23
%% TODO Anweisung: ein Satz über der ersten Teilaufgabe fehlt in der Bank – Nr. 23
\begin{aufgabe}{Zylinder}
%% TODO \rechenplatz{n}: Schrittzahl des Grundfalls fehlt in der Bank (nur bei fester Schreibform, 2.3 b)
\begin{teile}
\teil Eine Tasse ist oben $8$ cm breit, von Rand zu Rand durch die Mitte gemessen. Ist das Radius oder Durchmesser? Gib den Radius an. Radius \leerfeld[cm]
\teil Zylinder: Radius $5$ cm, Höhe $8$ cm – Volumen auf eine Stelle? \leerfeld[cm³]
\teil Zylinder: Durchmesser $8$ cm, Höhe $5$ cm – Volumen auf eine Stelle? \leerfeld[cm³]
\teil Zylinder: Radius $2{,}5$ cm, Höhe $7{,}5$ cm – Volumen auf eine Stelle? \leerfeld[cm³]
\teil Ein Kochtopf ist innen zylinderförmig: Radius $11$ cm, Höhe $14$ cm. Wie viele Liter passen hinein (auf eine Stelle)? \leerfeld[l]
\teil Zylinder: Radius $4$ cm, Höhe $9$ cm – Mantel auf eine Stelle? \leerfeld[cm²]
\teil Eine Dose mit Deckel: Radius $5$ cm, Höhe $9$ cm – Oberfläche auf eine Stelle? \leerfeld[cm²]
\teil Ein Zylinder hat den Radius $2$ cm und die Höhe $5$ cm. Welches Rechteck gehört zu seinem Netz? \\ \kreuz{$5$ cm × $6{,}3$ cm} \\ \kreuz{$5$ cm × $12{,}6$ cm} \\ \kreuz{$5$ cm × $4$ cm}
\teil Zylinder: Volumen $500$ cm³, Radius $5$ cm – Höhe auf eine Stelle? \leerfeld[cm]
\teil Ein größerer Becher soll bei $9$ cm Höhe $600$ cm³ fassen. Berechne den Radius auf eine Stelle \hfill (P10 2022 OS). \leerfeld[cm]
\teil Eine zylinderförmige Dose hat innen den Radius $3{,}5$ cm und die Höhe $9$ cm. Der Hersteller gibt „ca. 350 ml“ an. Weise nach, dass die Angabe stimmt \hfill (P10 2023 OS). Hinweis: $1$ cm³ $= 1$ ml.
\teil Zylinder A: Radius $3$ cm, Höhe $5$ cm. Zylinder B: Radius $6$ cm, Höhe $5$ cm. Wievielmal so groß ist das Volumen von B? Begründe mit der Rechnung. \leerfeld -mal so groß
\teil Eine Farbdose soll $2{,}5$ Liter fassen; ihr Deckel hat den Radius $7$ cm. Berechne die Innenhöhe der Dose auf eine Stelle \hfill (P10 2023 OS). \leerfeld[cm]
\end{teile}
\end{aufgabe}

%% TODO Titel: Ich-kann-Satz fehlt in der Bank (Platzhalter: Kettenname) – Nr. 24
%% TODO Anweisung: ein Satz über der ersten Teilaufgabe fehlt in der Bank – Nr. 24
\begin{aufgabe}{Zylinder – Fehler finden, Begründen, Anwendung, Darstellungswechsel}
\begin{teile}
\teil Eine Regentonne hat den Durchmesser $0{,}7$ m und die Höhe $0{,}9$ m. Max rechnet das Volumen: \rechnung{V &= \pi \cdot 0{,}7^2 \cdot 0{,}9 \\ V &\approx 1{,}39 \text{ m}^3} Wo ist der Fehler?
\teil Warum wird das Volumen eines Zylinders viermal so groß, wenn man den Radius verdoppelt?
\teil Eine Regentonne ist innen $56$ cm breit (Durchmesser) und $90$ cm hoch. An einem heißen Tag braucht der Garten $140$ l Wasser. Reicht eine volle Tonne für zwei heiße Tage?
\teil Skizziere das Netz eines Zylinders mit Radius $1{,}5$ cm und Höhe $4$ cm und schreibe die Maße des Rechtecks dazu.

\rechenplatz{5}
\end{teile}
\end{aufgabe}
